% Generated draft: requires mathematical and source review.
\documentclass[11pt]{article}
\usepackage[margin=25mm]{geometry}
\usepackage{fontspec}
\setmainfont{Latin Modern Roman}
\usepackage{amsmath,amssymb,mathtools,unicode-math}
\setmathfont{Latin Modern Math}
\usepackage{enumitem,needspace,xurl,hyperref}
\setlength{\parindent}{0pt}
\setlength{\parskip}{5pt}
\setlength{\emergencystretch}{3em}
\newcommand{\N}{\mathbb{N}}
\newcommand{\Z}{\mathbb{Z}}
\newcommand{\Q}{\mathbb{Q}}
\newcommand{\R}{\mathbb{R}}
\newcommand{\C}{\mathbb{C}}
\newcommand{\F}{\mathbb{F}}
\newcommand{\A}{\mathbb{A}}
\newcommand{\PP}{\mathbb{P}}
\newcommand{\E}{\mathbb{E}}
\newcommand{\eps}{\varepsilon}
\newcommand{\dd}{\,\mathrm{d}}
\begin{document}
\section*{UnsolvedMath ID: 700420}
\textbf{Source problem code:} catalogue entry
\subsection{Definitions and mathematical statement}
Let \(X\) be a finite set, let \(F(X)\) be the free group on \(X\), and let \(R\subseteq F(X)\) be a finite set of words. The finite presentation \(P=\langle X\mid R\rangle\) defines the group
\[
G_P=F(X)/\langle\!\langle R\rangle\!\rangle,
\]
where \(\langle\!\langle R\rangle\!\rangle\) is the normal closure of \(R\) in \(F(X)\). For a group \(G\), define \(G^{(0)}=G\) and
\[
G^{(n+1)}=[G^{(n)},G^{(n)}]\qquad(n\geq 0),
\]
where \([H,H]\) is the subgroup generated by all commutators \(a^{-1}b^{-1}ab\) with \(a,b\in H\). A group \(G\) is solvable if \(G^{(d)}=\{1\}\) for some integer \(d\geq 0\), and is metabelian if \(G^{(2)}=\{1\}\), equivalently if its commutator subgroup is abelian.

Does there exist an algorithm which, for every finite presentation \(P\) such that \(G_P\) is solvable, halts and outputs \(\mathrm{YES}\) if \(G_P^{(2)}=\{1\}\) and \(\mathrm{NO}\) otherwise?

Here solvability is a promise about the presented group: the input is its finite presentation, without an additional certificate of solvability or a bound on its derived length. The source does not specify behavior on presentations of nonsolvable groups; this formulation imposes no requirement on those inputs.
\subsection{Short English statement}
Given a finite presentation of a group promised to be solvable, can one always decide whether its commutator subgroup is abelian? This is the question attributed to F. B. Cannonito in Kourovka Notebook Problem 7.21.
\subsection{Sources}
\begin{itemize}
\item Local UnsolvedMath registry, permanent ID \(700420\), source problem code \texttt{catalogue entry}, titled \emph{Kourovka Notebook Problem 7.21}; supplied upstream revision \texttt{d371ee917ac9647fbb06cd6fc529deec6f3dda82}.
\item F. B. Cannonito, problem proposer, 1980, as attributed in the supplied record.
\item \emph{The Kourovka Notebook (Unsolved Problems in Group Theory)}, 7th issue (1980), Problem 7.21; the supplied record references the 21st issue edition (2026), \texttt{arXiv:1401.0300v48}. The collection description names E. I. Khukhro and V. D. Mazurov as editors.
\item The registry labels the problem \emph{open}. This is supplied source metadata, not an independently verified assessment of its mathematical status. The cited source and attribution have not been independently checked for this draft.
\end{itemize}
\end{document}
